
\chapter{\texorpdfstring{\gls{NAT} und Port Forwarding}{NAT und Port Forwarding}}
\label{app:nat_port_forwarding}
\section{\texorpdfstring{\acrfull{NAT}}{Network Address Translation (NAT)}}
\textbf{\gls{NAT}} (Network Address Translation) ist ein Verfahren, das in Routern verwendet wird, um die Kommunikation zwischen Computern in einem lokalen Netzwerk und dem Internet zu ermöglichen. Dabei wird die lokale \gls{IP}-Adresse eines Computers in eine globale \gls{IP}-Adresse übersetzt, wenn Datenpakete das lokale Netzwerk verlassen, und umgekehrt, wenn Datenpakete aus dem Internet empfangen werden. Dies ermöglicht es mehreren Computern in einem lokalen Netzwerk, dieselbe globale \gls{IP}-Adresse zu teilen, wodurch die Anzahl der benötigten globalen Adressen reduziert wird.

\begin{myexample}[\acrfull{NAT}]
Stellen Sie sich vor, Sie haben ein lokales Netzwerk mit mehreren Computern, die alle die gleiche globale \gls{IP}-Adresse \texttt{203.0.113.5} teilen. Ein Computer im lokalen Netzwerk mit der lokalen Adresse \texttt{192.168.1.10} möchte eine Anfrage an einen externen Server mit der Adresse \texttt{8.8.8.8} senden. Der Router übersetzt diese Anfrage wie folgt:

\begin{table}[H]
	\centering
	\begin{tabular}{c|c|c|c}
		\textbf{Phase} & \textbf{Quell-IP} & \textbf{Ziel-IP} & \textbf{Ort} \\
		\midrule\midrule
		Vor \gls{NAT} & \texttt{192.168.1.10} (lokal) & \texttt{8.8.8.8} & Lokales Netzwerk \\
		\midrule
		Nach \gls{NAT} & \texttt{203.0.113.5} (global) & \texttt{8.8.8.8} & Internet \\
		\midrule
		Rückweg & \texttt{8.8.8.8} & \texttt{203.0.113.5} (global) & Internet \\
		\midrule
		Nach Rückübersetzung & \texttt{8.8.8.8} & \texttt{192.168.1.10} (lokal) & Lokales Netzwerk \\
	\end{tabular}
	\caption{Beispiel einer \gls{NAT}-Übersetzung: Ein lokales Gerät kommuniziert mit einem externen Server}
\end{table}

Dadurch können mehrere Geräte im lokalen Netzwerk gleichzeitig mit verschiedenen (oder sogar denselben) externen Servern kommunizieren, obwohl sie alle die gleiche globale \gls{IP}-Adresse nutzen (und potentiell sogar über denselben Port kommunizieren). Der Router kann anhand der durch ihn zugewiesenen Port-Nummern und lokalen Adressen genau unterscheiden, welche Antworten an welche lokalen Geräte zurückgeleitet werden müssen.
\end{myexample}

\gls{NAT} wird weiter unterschieden in:
\begin{itemize}
    \item \textbf{Statisches NAT}: Hierbei wird eine feste Zuordnung zwischen jeder lokalen \gls{IP}-Adresse und einer globalen \gls{IP}-Adresse erstellt. Dies wird oft für Server verwendet, die von aussen erreichbar sein müssen. Für private Netzwerke mit vielen Geräten ist dies jedoch unpraktisch, da es viele globale \gls{IP}-Adressen erfordert.
    \item \textbf{Dynamisches NAT}: Hierbei wird eine lokale \gls{IP}-Adresse dynamisch einer verfügbaren globalen \gls{IP}-Adresse zugeordnet, wenn eine Verbindung nach aussen hergestellt wird. Dies ist effizienter als statisches \gls{NAT}, da nicht für jedes Gerät eine feste globale Adresse benötigt wird. Aufgrund der Tatsache, dass mehrere \gls{IP}-Adressen benötigt werden, ist dies jedoch immer noch nicht ideal für grosse Netzwerke.
    \item \textbf{\gls{PAT} (Port Address Translation)}: Auch als \textbf{NAT Overloading} bekannt, ermöglicht \gls{PAT} mehreren Geräten in einem lokalen Netzwerk, dieselbe globale \gls{IP}-Adresse zu teilen, indem unterschiedliche Port-Nummern verwendet werden. Dies ist die am häufigsten verwendete Form von \gls{NAT} in Heimnetzwerken, da es die Anzahl der benötigten globalen \gls{IP}-Adressen minimiert. Jeder lokale Computer wird durch eine Kombination aus der gemeinsamen globalen \gls{IP}-Adresse und einer eindeutigen Port-Nummer identifiziert In Praxis wird mit \gls{NAT} indirekt meistens \gls{PAT} gemeint.
\end{itemize}

\section{\texorpdfstring{\acrfull{PAT}}{Port Address Translation (PAT)}}
Angenommen, Sie haben drei Computer in Ihrem lokalen Netzwerk mit den folgenden lokalen \gls{IP}-Adressen:
\begin{itemize}
    \item Computer 1: \texttt{192.168.1.10}
    \item Computer 2: \texttt{192.168.1.11}
    \item Computer 3: \texttt{192.168.1.12}
\end{itemize}
Alle drei Computer teilen sich die gleiche globale \gls{IP}-Adresse \texttt{203.0.113.5}. Wenn alle drei Computer gleichzeitig eine Verbindung zu einem externen Server herstellen möchten, verwendet der Router nicht nur die \gls{IP}-Adressen zur Unterscheidung, sondern auch die \textbf{Port-Nummern}. Der Router erstellt dabei eine \textbf{\gls{PAT}-Tabelle} (häufig auch \textbf{\gls{NAT}-Tabelle} genannt), die nicht nur die lokalen und globalen \gls{IP}-Adressen, sondern auch die zugehörigen Ports speichert. Die Port-Nummern ermöglichen es dem Router, die eingehenden Antworten vom externen Server korrekt an die jeweiligen lokalen Computer weiterzuleiten, obwohl sie alle dieselbe globale \gls{IP}-Adresse verwenden.

Dabei ist es sogar möglich, dass mehrere lokale Geräte über denselben lokalen Port kommunizieren, solange die Kombination aus globaler \gls{IP}-Adresse und globalem Port eindeutig bleibt. Die \gls{PAT}-Tabelle im Router könnte wie folgt aussehen:

\begin{table}[H]
    \centering
    \begin{tabular}{c|c|c|c|c}
        \textbf{Lokale \gls{IP}:Port} & \textbf{Globale \gls{IP}:Port} & \textbf{Ziel-\gls{IP}:Port} & \textbf{Protokoll} & \textbf{Status} \\
        \midrule\midrule
        \texttt{192.168.1.10:45000} & \texttt{203.0.113.5:45000} & \texttt{8.8.8.8:443} & \gls{TCP} & Aktiv \\
        \midrule
        \texttt{192.168.1.11:45000} & \texttt{203.0.113.5:45001} & \texttt{8.8.8.8:443} & \gls{TCP} & Aktiv \\
        \midrule
        \texttt{192.168.1.12:45000} & \texttt{203.0.113.5:45002} & \texttt{93.184.216.34:80} & \gls{TCP} & Aktiv \\
    \end{tabular}
    \caption{\gls{PAT}-Tabelle mit mehreren lokalen Geräten, die gleichzeitig auf externe Server zugreifen}
\end{table}

In diesem Beispiel verwenden alle drei lokalen Computer zufälligerweise denselben lokalen Port (45000), erhalten jedoch vom Router unterschiedliche globale Ports (45000, 45001, 45002) zugewiesen. Zwei Computer (192.168.1.10 und 192.168.1.11) greifen auf denselben externen Server \texttt{8.8.8.8:443} zu, während der dritte Computer (192.168.1.12) auf einen anderen Server \texttt{93.184.216.34:80} zugreift. Der Router kann die eingehenden Antworten anhand der Kombination aus globaler \gls{IP}-Adresse und globalem Port eindeutig dem richtigen lokalen Computer zuordnen.

\section{Port Forwarding}
Port Forwarding kann verwendet werden, um eingehende Verbindungen von externen Computern an bestimmte lokale Computer weiterzuleiten. Dies ist besonders nützlich, wenn ein lokaler Computer (z.B. ein Webserver oder ein Spieleserver) von ausserhalb des lokalen Netzwerks erreichbar sein soll. Durch das Einrichten von Port Forwarding im Router kann der Datenverkehr, der an eine bestimmte Port-Nummer der globalen \gls{IP}-Adresse gesendet wird, an die entsprechende lokale \gls{IP}-Adresse und Port-Nummer weitergeleitet werden. 

\begin{myexample}[Port Forwarding]
Angenommen, Sie haben einen Webserver in Ihrem lokalen Netzwerk mit der lokalen \gls{IP}-Adresse \texttt{192.168.1.20} und möchten, dass dieser von ausserhalb Ihres Netzwerks über das Internet erreichbar ist. Ihr Router hat die globale \gls{IP}-Adresse \texttt{203.0.113.5}. Sie richten im Router eine Port-Weiterleitung ein, die eingehende Verbindungen auf Port 80 (\gls{HTTP}) der globalen Adresse an die lokale Adresse \texttt{192.168.1.20} auf Port 80 weiterleitet. Dadurch können externe Benutzer über \texttt{http://203.0.113.5} auf Ihren Webserver zugreifen, obwohl dieser sich in einem privaten Netzwerk befindet. Die Forwarding-Tabelle im Router könnte wie folgt aussehen:
\begin{table}[H]
	\centering
	\begin{tabular}{c|c|c|c}
		\textbf{Externer Port} & \textbf{Interne IP} & \textbf{Interner Port} & \textbf{Protokoll} \\
		\midrule\midrule
		80 & \texttt{192.168.1.20} & 80 & \gls{TCP} \\
	\end{tabular}
	\caption{Port Forwarding-Tabelle im Router}
\end{table}

Das Protokoll kann entweder \gls{TCP} oder \gls{UDP} sein, je nachdem, welche Art von Datenverkehr weitergeleitet werden soll. Typischerweise wird \gls{TCP} für Webserver und andere verbindungsorientierte Dienste verwendet, da es eine zuverlässige Datenübertragung gewährleistet. \gls{UDP} wird hingegen typischerweise für Anwendungen wie Online-Spiele oder \gls{VPN}-Dienste genutzt, die eine schnellere, aber weniger zuverlässige Übertragung erfordern.
\end{myexample}

\begin{myexample}[\acrshort{NAS} im lokalen Netzwerk von aussen erreichbar machen]
	Angenommen, Sie kaufen ein \gls{NAS} für Ihr Heimnetzwerk, um Dateien zentral zu speichern und von verschiedenen Geräten darauf zuzugreifen. Ihr \gls{NAS} hat die lokale \gls{IP}-Adresse \texttt{192.168.1.20}. Sie möchten jedoch auch von ausserhalb Ihres Heimnetzwerks auf das \gls{NAS} zugreifen können, z.B. wenn Sie unterwegs sind. Ihr Internetanbieter hat Ihnen die globale \gls{IP}-Adresse \texttt{203.0.113.5} zugewiesen. Zudem haben Sie sich eine statische \gls{IP}v4-Adresse gekauft, damit sich Ihre \gls{IP}-Adresse nicht ändert.

	Um dies zu ermöglichen, müssen Sie Port Forwarding auf Ihrem Router einrichten. Angenommen, Ihr \gls{NAS} verwendet den Standardport 8080 für den Zugriff auf die Weboberfläche. Sie richten im Router eine Port-Weiterleitung ein, die eingehende Verbindungen auf den Port 8080 der globalen Adresse \texttt{203.0.113.5} an die lokale Adresse \texttt{192.168.1.20} weiterleitet. Dadurch können Sie von ausserhalb Ihres Heimnetzwerks auf Ihr \gls{NAS} zugreifen, indem Sie in Ihrem Webbrowser die Adresse \texttt{http://203.0.113.5:8080} eingeben.
	
	Die Forwarding-Tabelle im Router könnte wie folgt aussehen:
	\begin{table}[H]
		\centering
		\begin{tabular}{c|c|c|c}
			\textbf{Externer Port} & \textbf{Interne IP} & \textbf{Interner Port} & \textbf{Protokoll} \\
			\midrule\midrule
			8080 & \texttt{192.168.1.20} & 8080 & \gls{TCP} \\
	\end{tabular}
	\caption{Port Forwarding-Tabelle im Router}
\end{table}

	Damit Ihr \gls{NAS} auch wirklich immer die lokale \gls{IP}-Adresse \texttt{192.168.1.20} behält, sollten Sie im Router eine \gls{IP}-Reservierung für die \gls{MAC}-Adresse Ihres \gls{NAS} einrichten (s. \cref{remark:ip-reservation}).
\end{myexample}

\begin{myremark}[Unterschied \gls{NAT} und Port Forwarding]
\gls{NAT} und Port Forwarding sind eng verwandte Konzepte, die häufig zusammen verwendet werden, aber unterschiedliche Funktionen erfüllen:

\begin{itemize}
	\item \textbf{\gls{NAT}}: Ist ein automatischer Prozess, der auf dem Router läuft und dafür sorgt, dass Computer im lokalen Netzwerk mit dem Internet kommunizieren können. \gls{NAT} übersetzt automatisch lokale \gls{IP}-Adressen in globale Adressen für ausgehende Verbindungen und führt die Rückübersetzung für eingehende Antworten durch. Dies geschieht transparent und erfordert keine manuelle Konfiguration pro Gerät.
	
	\item \textbf{Port Forwarding}: Ist eine manuell konfigurierte Funktionalität, die es ermöglicht, eingehende Verbindungen von aussen gezielt an spezifische lokale Computer weiterzuleiten. Während \gls{NAT} primär für ausgehende Verbindungen zuständig ist, ermöglicht Port Forwarding, dass externe Benutzer von aussen initiierte Verbindungen zu Servern im lokalen Netzwerk herstellen können.
	
	\item \textbf{Zusammenhang}: Port Forwarding funktioniert auf Basis von \gls{NAT}. Es nutzt die \gls{NAT}-Tabelle des Routers, um eingehende Verbindungen korrekt an die richtigen lokalen Geräte weiterzuleiten. Ohne \gls{NAT} wäre Port Forwarding nicht möglich.
\end{itemize}

Zusammengefasst: \gls{NAT} ermöglicht Computern im lokalen Netzwerk, nach aussen zu kommunizieren, während Port Forwarding es externen Computern ermöglicht, gezielt auf Dienste im lokalen Netzwerk zuzugreifen.
\end{myremark}